% =====================================================================
% Guide to the 1D-QM-Playground B-spline TISE solver.
%
% Build:  pdflatex student-guide.tex   (run twice for the table of contents)
%
% Every numerical result quoted in this document was produced by the code
% described in it. Where a table reports a difference, it is against a
% closed-form solution.
% =====================================================================
\documentclass[11pt]{article}

\usepackage[margin=1in]{geometry}
\usepackage{amsmath,amssymb}
\usepackage{listings}
\usepackage{xcolor}
\usepackage{booktabs}
\usepackage{enumitem}
\usepackage{microtype}
\usepackage[colorlinks=true,linkcolor=blue!50!black,urlcolor=blue!50!black]{hyperref}

\definecolor{codebg}{gray}{0.96}
\definecolor{keyw}{rgb}{0.1,0.3,0.6}
\definecolor{comm}{rgb}{0.3,0.5,0.3}
\definecolor{strc}{rgb}{0.6,0.2,0.2}

\lstdefinestyle{base}{
  backgroundcolor=\color{codebg},
  basicstyle=\ttfamily\small,
  commentstyle=\color{comm}\itshape,
  keywordstyle=\color{keyw}\bfseries,
  stringstyle=\color{strc},
  breaklines=true,
  breakatwhitespace=false,
  columns=fullflexible,
  frame=single,
  rulecolor=\color{gray!40},
  framesep=4pt,
  xleftmargin=6pt,
  showstringspaces=false,
  captionpos=b,
}
\lstdefinelanguage{yamlx}{
  keywords={run,physics,bspline,potential,potential_deltas,tise,continuum,
            visualization,run_tise,run_tdse,run_analysis,output_dir,mass,hbar,
            n_nodes,order,domain,n_pts_eigenstate,error_threshold,enabled,
            E_threshold,E_max,n_energies,n_pts,l,eigenstates,phase_shifts,
            spectrum_overview,x,strength,true,false},
  sensitive=true,
  comment=[l]{\#},
  morestring=[b]',
}
\lstdefinestyle{yaml}{style=base,language=yamlx}
\lstdefinestyle{py}{style=base,language=Python}
\lstdefinestyle{sh}{style=base,language=bash}

% Long unhyphenatable \texttt paths (e.g. examples/confining-well.yaml) would
% otherwise push past the right margin; this lets TeX stretch a line rather
% than overflow it.
\setlength{\emergencystretch}{3em}

\newcommand{\file}[1]{\texttt{#1}}
\newcommand{\bra}[1]{\langle #1 |}
\newcommand{\ket}[1]{| #1 \rangle}
\newcommand{\braket}[2]{\langle #1 | #2 \rangle}

\title{\textbf{The B-Spline TISE Solver}\\[2ex]
       \large Computing bound and continuum states of a one-dimensional potential}
\author{}
\date{}

\begin{document}
\maketitle
\thispagestyle{empty}

\begin{center}
\begin{minipage}{0.88\textwidth}
\small
This is a guide to the program: how to build it, how to describe a potential to
it, how to read what it produces, and how to choose the parameters that control
accuracy. It assumes you know the quantum mechanics and have never seen this
codebase. Sections~\ref{sec:build}--\ref{sec:output} take about twenty minutes
and end with a solved problem and a plot.
\end{minipage}
\end{center}

\vspace{2ex}
\tableofcontents
\newpage

% =====================================================================
\section{What the program computes}
\label{sec:scope}

Given a potential $V(x)$ that you write in a configuration file, the solver
finds the eigenstates of
\begin{equation}
  -\frac{\hbar^{2}}{2m}\,\psi''(x) + V(x)\,\psi(x) = E\,\psi(x),
  \qquad x \in [x_{\min}, x_{\max}].
  \label{eq:tise}
\end{equation}
It expands $\psi$ in a B-spline basis $\{B_j\}$, $\psi(x) = \sum_j c_j B_j(x)$,
turning~\eqref{eq:tise} into the generalized matrix eigenproblem
\begin{equation}
  H\mathbf{c} = E\,S\,\mathbf{c},
  \qquad
  H_{ij} = \frac{\hbar^{2}}{2m}\!\int\! B_i' B_j'\,dx + \!\int\! B_i V B_j\,dx,
  \qquad
  S_{ij} = \int\! B_i B_j\,dx .
  \label{eq:gep}
\end{equation}

You get every eigenvalue $E_n$ and eigenfunction $\psi_n(x)$, tabulated on a
grid of your choosing, and --- on request --- continuum (scattering) states
with their phase shifts $\delta(E)$. The potential may be piecewise,
discontinuous, singular, or carry Dirac $\delta$ terms; the basis adapts to
each of these on its own (Sections~\ref{sec:structure} and~\ref{sec:delta}).

The wavefunction vanishes at both ends of the box,
$\psi(x_{\min}) = \psi(x_{\max}) = 0$. Place the walls beyond the point where
the states you care about have decayed, and they will not affect your answer;
Section~\ref{sec:convergence} makes that quantitative.

Time evolution is done from these eigenstates, by giving each one its phase.
That calculation is about ten lines of Python and is the subject of
Section~\ref{sec:evolution}.

% =====================================================================
\section{Installing and building}
\label{sec:build}

You need a C++17 compiler, CMake $\geq 3.20$, Git, and Python~3.9 or newer.
Every C++ library the solver uses is located on your system or, failing that,
downloaded and compiled for you during the CMake configure step. The commands
below are the same on macOS, Linux and Windows.

\begin{lstlisting}[style=sh]
git clone -b PHY5606_Projects https://github.com/Razgriz117/BSpline-CPP.git
cd BSpline-CPP/1D-QM-Playground

cmake -S TISE -B TISE/build -DCMAKE_BUILD_TYPE=Release -DBUILD_TESTING=ON
cmake --build TISE/build --config Release

pip install -r requirements.txt
\end{lstlisting}

The first configure fetches dependencies, so it needs a network connection and
takes a couple of minutes; later builds reuse them.

Pass both \texttt{-DCMAKE\_BUILD\_TYPE=Release} and \texttt{--config Release},
as above. Single-config generators (Unix Makefiles, Ninja) read the first;
multi-config generators (Visual Studio, Xcode) read the second. Giving both
produces an optimized build everywhere.

On Windows, run them in a ``Developer Command Prompt for VS'', or any shell
with MSVC on the \texttt{PATH}. The executable appears at
\file{TISE/build/Release/tise\_solver.exe} and the Python front end finds it
for you. Building inside WSL works too, and is then the Linux case.

Confirm the installation:
\begin{lstlisting}[style=sh]
ctest --test-dir TISE/build --build-config Release
python3 controller.py --config examples/confining-well.yaml
\end{lstlisting}
The first reports that all test suites pass. The second solves a finite square
well and writes results into \file{data/confining-well/}.

% =====================================================================
\section{Running a calculation}

Drive everything through the Python front end:
\begin{lstlisting}[style=sh]
python3 controller.py --config my-problem.yaml
\end{lstlisting}
which runs the solver and then \file{analysis.py} to produce plots. For the
data files alone, call the binary directly:
\begin{lstlisting}[style=sh]
./TISE/build/tise_solver --config my-problem.yaml --output-dir out/
\end{lstlisting}

A calculation is described entirely by its configuration file. Copy
\file{examples/confining-well.yaml} and edit it.

% =====================================================================
\section{The input file}
\label{sec:input}

\begin{lstlisting}[style=yaml,caption={A complete input file.}]
run:
  run_tise:     true
  run_tdse:     false
  run_analysis: true         # also produce plots
  output_dir:   "./data/my-problem"

physics:
  mass: 1.0                  # m in equation (1)
  hbar: 1.0                  # atomic units

bspline:
  n_nodes: 201               # resolution
  order:   12                # B-spline order k
  domain:  [-10.0, 10.0]     # the box [x_min, x_max]

potential:
  - "{'domain': '[-10, 10]', 'function': '0.5*x^2'}"

# potential_deltas:          # Dirac delta terms; see section 10
#   - {x: 0.0, strength: -1.0}

tise:
  n_pts_eigenstate: 2001     # points at which each psi_n is tabulated
  error_threshold:  1.0e-10
  continuum:
    enabled: false

visualization:
  spectrum_overview: true    # V(x) with each psi_n drawn on its level E_n
  eigenstates:       true
\end{lstlisting}

\subsection{Describing the potential}

\texttt{potential} is a list of pieces that together tile
\texttt{bspline.domain}. Each piece is a string holding a small dictionary:

\begin{itemize}[nosep]
\item \texttt{domain} --- interval notation. Square brackets are closed, round
      brackets open, so \texttt{'[0, 5)'} means $0 \le x < 5$. \texttt{inf} is
      available, as in \texttt{'(0, inf)'}.
\item \texttt{function} --- an expression in \texttt{x}, evaluated at run time.
      You have \texttt{+ - * / \^{}}, \texttt{sqrt}, \texttt{exp},
      \texttt{log}, \texttt{sin}, \texttt{cos}, \texttt{tan}, \texttt{abs} and
      the usual companions.
\end{itemize}

A square well with a wall at the origin, in two pieces:
\begin{lstlisting}[style=yaml]
potential:
  - "{'domain': '[0, 10)',   'function': '-1.0'}"
  - "{'domain': '[10, 100]', 'function': '0'}"
\end{lstlisting}

Write every constant in \texttt{function} as a literal number. The expression
is evaluated on its own, so \texttt{mass}, \texttt{hbar} and the other
configuration fields are not in scope: for a mass-dependent coefficient, work
the number out and write it in. When you change \texttt{physics.mass}, revisit
any constant in \texttt{function} that depended on it.

For a fixed-$\ell$ radial problem, include the centrifugal term in the
expression. Hydrogen's $V_\ell(r) = \ell(\ell+1)/2r^{2} - 1/r$ at $\ell = 1$
becomes \texttt{'-1/x + 1/x\^{}2'}.

Dirac $\delta$ terms have their own \texttt{potential\_deltas} section, since a
$\delta$ is not an expression in $x$ (Section~\ref{sec:delta}). They combine
freely with whatever smooth \texttt{potential} you write.

\subsection{Scanning a parameter}

To vary a parameter, generate the configuration files:

\begin{lstlisting}[style=py]
import subprocess, yaml, pathlib

base = yaml.safe_load(open("examples/confining-well.yaml"))
for depth in [0.5, 1.0, 2.0, 4.0]:
    cfg = dict(base)
    cfg["potential"] = [
        "{'domain': '[0, 10)',   'function': '%g'}" % (-depth),
        "{'domain': '[10, 100]', 'function': '0'}",
    ]
    cfg["run"]["output_dir"] = f"./data/scan/depth{depth}"
    path = pathlib.Path(f"cfg_depth{depth}.yaml")
    path.write_text(yaml.safe_dump(cfg))
    subprocess.run(["python3", "controller.py", "--config", str(path)], check=True)
\end{lstlisting}

% =====================================================================
\section{Reading the output}
\label{sec:output}

Results land in \file{<output\_dir>/tise/}.

\begin{center}
\begin{tabular}{@{}lp{0.62\textwidth}@{}}
\toprule
File & Contents \\
\midrule
\file{eigenvalues.dat}     & two columns: index (0-based), $E_n$ \\
\file{eigenstates.dat}     & column 1 is $x$; column $n{+}2$ is $\psi_n(x)$. All
                             eigenvalues are repeated in the \texttt{\#} header \\
\file{continuum\_states.dat} & the same layout, for scattering states \\
\file{phase\_shifts.dat}   & three columns: $E$, $\delta(E)$, $d\delta/dE$ \\
\file{potential.dat}       & $x$, $V(x)$ on the same grid \\
\file{warnings.json}       & the solver's physics notes on the run \\
\bottomrule
\end{tabular}
\end{center}

Loading them takes one line each:
\begin{lstlisting}[style=py]
import numpy as np
d   = np.loadtxt("data/my-problem/tise/eigenstates.dat")
E   = np.loadtxt("data/my-problem/tise/eigenvalues.dat")[:, 1]
x   = d[:, 0]
phi = d[:, 1:]            # phi[:, n] is psi_n(x)
\end{lstlisting}

The tabulated $\psi_n$ come out normalized: $\int |\psi_n|^{2} dx = 1$ and
$\int \psi_m \psi_n\,dx = \delta_{mn}$ to machine precision on the output grid.
Use them in quadrature directly.

\file{eigenvalues.dat} lists every state the basis produced. The bound states
are those with $E < V(\infty)$; the ones above that threshold are the finite
box's discretization of the continuum, and Section~\ref{sec:continuum} is how
to treat scattering properly.

The per-state files \file{eigenstate\_NNN.dat} are numbered from~1, while the
index column of \file{eigenvalues.dat} starts at~0. Working from
\file{eigenstates.dat} keeps one convention throughout.

Read \file{warnings.json} after each run. It carries the solver's physics
observations --- a state reaching the outer wall, a continuum energy close to a
box state, a basis coarser than the energy range requested --- and each one
points at a parameter worth adjusting.

% =====================================================================
\section{Choosing the basis and the box}
\label{sec:convergence}

Three parameters control accuracy.

\begin{description}[style=nextline,leftmargin=1em]
\item[{\texttt{bspline.domain} --- the box $[x_{\min}, x_{\max}]$.}]
  Place the walls beyond the classical turning point of the highest state you
  need, so $\psi$ has decayed before reaching them.
\item[\texttt{bspline.n\_nodes} --- resolution.] The principal knob. Double it
  and re-run to confirm the answer has settled.
\item[\texttt{bspline.order} --- the B-spline order $k$.] Each basis function is
  a piecewise polynomial of degree $k-1$. Use $8$--$12$. Higher order converges
  faster per node and widens the matrix band.
\end{description}

\subsection{Box size}

Hydrogen makes the requirement concrete. Holding \texttt{n\_nodes}${}=201$ fixed
and enlarging only the box, against $E_n = -1/2n^{2}$:

\begin{center}
\begin{tabular}{@{}rrrrr@{}}
\toprule
$R = x_{\max}$ & $n=1$ & $n=2$ & $n=3$ & $n=4$ \\
\midrule
20  & $3\times10^{-14}$ & $1\times10^{-5}$  & $6\times10^{-3}$  & $5\times10^{-2}$ \\
40  & $2\times10^{-15}$ & $5\times10^{-13}$ & $1\times10^{-6}$  & $7\times10^{-4}$ \\
80  & $6\times10^{-15}$ & $7\times10^{-15}$ & $2\times10^{-15}$ & $1\times10^{-9}$ \\
160 & $8\times10^{-15}$ & $9\times10^{-16}$ & $2\times10^{-15}$ & $2\times10^{-16}$ \\
\bottomrule
\end{tabular}
\end{center}

The pattern follows the classical turning point: hydrogenic state $n$ extends
to $r \approx 2n^{2}$, so $n=4$ reaches $r \approx 32$ and only starts to decay
beyond it. Choose $R$ from the highest state you intend to use --- for four
hydrogen states, $R \approx 160$ delivers all of them to machine precision.

A state that reaches the wall is reported in \file{warnings.json}:
\begin{lstlisting}[style=sh]
bound state 6 (E=-0.00965) appears to be colliding with the outer wall at
x=100 (psi'(rMax)=0.0104, exceeds tolerance) -- its energy and wavefunction
may be inaccurate due to box truncation.
\end{lstlisting}
Enlarging the domain clears it.

\subsection{Resolution}

Bound-state energies converge quickly. Phase shifts ask for more, because at
higher $E$ the wavefunction oscillates faster and the node spacing has to
resolve it. For a square well of depth~1 and width~10 on $[0,100]$ at
order~12, the difference in $\delta(E)$ from the closed form:

\begin{center}
\begin{tabular}{@{}rrrr@{}}
\toprule
\texttt{n\_nodes} & $E=0.1$ & $E=0.3$ & $E=0.5$ \\
\midrule
51  & $3\times10^{-6}$  & $2\times10^{-5}$  & $2\times10^{-4}$ \\
101 & $2\times10^{-11}$ & $2\times10^{-12}$ & $5\times10^{-9}$ \\
201 & $1\times10^{-13}$ & $2\times10^{-13}$ & $6\times10^{-13}$ \\
401 & $6\times10^{-13}$ & $4\times10^{-13}$ & $1\times10^{-12}$ \\
\bottomrule
\end{tabular}
\end{center}

At 51 nodes the bound states already reach $10^{-13}$ while $\delta(E{=}0.5)$
reaches $10^{-4}$; 201 nodes brings the whole range to $10^{-13}$. When you are
working with the continuum, take $\delta(E)$ as the quantity to converge. The
solver estimates the resolvable energy ceiling from your node spacing and
records it in \file{warnings.json}.

\subsection{\texttt{n\_pts\_eigenstate}}

This sets how finely the resulting $\psi_n$ are tabulated, independently of the
solve. Everything in Section~\ref{sec:evolution} is quadrature on that grid, so
set it generously --- 2001 is a good default.

% =====================================================================
\section{Evolving a state in time}
\label{sec:evolution}

A state built from eigenstates of a time-independent Hamiltonian evolves by
phases alone. Writing
\begin{equation}
  \ket{\psi(0)} = \sum_n c_n \ket{\phi_n},
  \qquad c_n = \braket{\phi_n}{\psi(0)},
\end{equation}
and using $\hat H \ket{\phi_n} = E_n \ket{\phi_n}$,
\begin{equation}
  \boxed{\;
  \psi(x,t) = \sum_n c_n\, \phi_n(x)\, e^{-i E_n t/\hbar}.
  \;}
  \label{eq:spectral}
\end{equation}
This is exact: there is no time step to choose and no propagation error to
control. The solver supplies $\phi_n(x)$ in \file{eigenstates.dat} and $E_n$ in
\file{eigenvalues.dat}; the rest is two integrals and a matrix product.

\subsection{The recipe}

The tabulated eigenstates are orthonormal on the output grid, so the overlap
integrals are ordinary quadrature.

\begin{lstlisting}[style=py,caption={A complete spectral-propagation helper.}]
import numpy as np

def load(run_dir):
    """Return x, phi (columns are phi_n), and E."""
    d = np.loadtxt(f"{run_dir}/eigenstates.dat")
    E = np.loadtxt(f"{run_dir}/eigenvalues.dat")[:, 1]
    return d[:, 0], d[:, 1:], E

def coefficients(x, phi, psi0, n_states):
    """c_n = <phi_n | psi(0)>, by quadrature on the output grid."""
    return np.array([np.trapezoid(phi[:, n].conj() * psi0, x)
                     for n in range(n_states)])

def propagate(x, phi, E, c, t, hbar=1.0):
    """psi(x, t) = sum_n c_n phi_n(x) exp(-i E_n t / hbar)."""
    n = len(c)
    return phi[:, :n] @ (c * np.exp(-1j * E[:n] * t / hbar))
\end{lstlisting}

\texttt{np.trapezoid} arrived in NumPy~2.0; on NumPy~1.x the same function is
\texttt{np.trapz}. To work with either, begin with
\begin{lstlisting}[style=py]
trapezoid = getattr(np, "trapezoid", None) or np.trapz
\end{lstlisting}
and call \texttt{trapezoid(...)} throughout.

Choose the number of states $N$ so that $\sum_{n<N} |c_n|^{2}$ reaches 1 to the
precision you want, using energies that are themselves converged
(Section~\ref{sec:convergence}). Printing that sum tells you at a glance
whether the basis represents your initial state.

The same \texttt{coefficients()} call gives the overlap of the eigenstates with
any function you supply, which is how you project an arbitrary wavepacket ---
a shifted Gaussian, say --- onto the spectrum.

\subsection{Checking it}

A coherent state in a harmonic well gives an exact benchmark, since it follows
the classical trajectory. With $V = x^{2}/2$, $m=1$ on $[-10,10]$ and
$\psi(x,0) = \pi^{-1/4} e^{ik_0 x} e^{-(x-x_0)^{2}/2}$ at $x_0=2$, $k_0=1$,
theory gives $\langle x\rangle(t) = x_0\cos t + k_0 \sin t$. Using $N=40$
states and the code above:

\begin{center}
\begin{tabular}{@{}rrrr@{}}
\toprule
$t$ & $\langle x\rangle$ numerical & analytic & difference \\
\midrule
$0$       & $2.000000000$  & $2.000000000$  & $4\times10^{-16}$ \\
$0.5$     & $2.234590662$  & $2.234590662$  & $4\times10^{-15}$ \\
$\pi/2$   & $1.000000000$  & $1.000000000$  & $1\times10^{-14}$ \\
$\pi$     & $-2.000000000$ & $-2.000000000$ & $2\times10^{-14}$ \\
$2\pi$    & $2.000000000$  & $2.000000000$  & $7\times10^{-15}$ \\
$10$      & $-2.222164169$ & $-2.222164169$ & $1\times10^{-14}$ \\
\bottomrule
\end{tabular}
\end{center}

with $\sum_n |c_n|^{2} = 1.000000000000$.

\subsection{Quantities computed from \texorpdfstring{$\psi(x,t)$}{psi(x,t)}}

\begin{lstlisting}[style=py]
rho  = np.abs(psi_t)**2
norm = np.trapezoid(rho, x)                       # stays 1
xbar = np.trapezoid(rho * x, x) / norm            # <x>(t)
P_ab = np.trapezoid(rho[(x>=a)&(x<=b)], x[(x>=a)&(x<=b)])  # P_[a,b](t)
pops = np.abs(c)**2                               # p_n, constant in time
surv = np.abs(np.sum(np.abs(c)**2 * np.exp(-1j*E[:len(c)]*t)))**2  # |<psi(0)|psi(t)>|^2
\end{lstlisting}

The survival probability $|\braket{\psi(0)}{\psi(t)}|^{2}$ is the natural
observable for tunneling and decay times.

% =====================================================================
\section{Worked potentials}
\label{sec:examples}

Five potentials that exercise different parts of the input language, with the
accuracy each reaches.

\subsection*{Harmonic oscillator}
\begin{lstlisting}[style=yaml]
physics: {mass: 1.0, hbar: 1.0}
bspline: {n_nodes: 151, order: 12, domain: [-10.0, 10.0]}
potential:
  - "{'domain': '[-10, 10]', 'function': '0.5*x^2'}"
\end{lstlisting}
Gives $E_n = n + \tfrac12$ to $10^{-14}$. Changing \texttt{mass} to 100 gives
$E_n = 0.1(n+\tfrac12)$ to the same accuracy, since $\omega = \sqrt{k/m}$ and
the kinetic term carries the mass.

\subsection*{Quartic double well}
\begin{lstlisting}[style=yaml]
physics: {mass: 1.0, hbar: 1.0}
bspline: {n_nodes: 201, order: 12, domain: [-10.0, 10.0]}
potential:
  - "{'domain': '[-10, 10]', 'function': '(1/200)*x^4 - 0.25*x^2 + 1/(64*(1/200))'}"
\end{lstlisting}
With $V = \alpha x^{4} - x^{2}/4 + 1/64\alpha$ at $\alpha = 1/200$, the minima
sit at $x = \pm5$ with $V=0$ and the central barrier is $V(0) = 3.125$. The
tunneling doublets come out directly:
\begin{center}
\begin{tabular}{@{}lrrr@{}}
\toprule
doublet & $E$ (lower) & $E$ (upper) & splitting \\
\midrule
1 & $0.48949752$ & $0.48949813$ & $6.1\times10^{-7}$ \\
2 & $1.42183914$ & $1.42193174$ & $9.3\times10^{-5}$ \\
3 & $2.26552000$ & $2.27064616$ & $5.1\times10^{-3}$ \\
\bottomrule
\end{tabular}
\end{center}
The splitting grows by roughly two orders of magnitude per doublet as the
states approach the barrier top, and the tunneling time follows as
$T = 2\pi\hbar/\Delta E$. A superposition of a doublet, localized in one well,
is built by taking $(\phi_1 \pm \phi_2)/\sqrt{2}$ with the sign that matches
the parity of the states you obtained.

\subsection*{Morse potential}
\begin{lstlisting}[style=yaml]
physics: {mass: 100.0, hbar: 1.0}
bspline: {n_nodes: 201, order: 12, domain: [0.0, 10.0]}
potential:
  - "{'domain': '(0, 10]', 'function': '(1-exp(-(x-1)))^2'}"
\end{lstlisting}
A hard wall at $x=0$ needs nothing beyond starting the domain there, since
$\psi(x_{\min})=0$ already. Against the analytic Morse spectrum
\begin{equation*}
  E_n = \omega_0\left(n+\tfrac12\right)
      - \frac{\left[\omega_0\left(n+\tfrac12\right)\right]^{2}}{4D_e},
  \qquad \omega_0 = a\sqrt{2D_e/m},
\end{equation*}
with $D_e=1$, $a=1$, $x_e=1$, $m=100$:
\begin{center}
\begin{tabular}{@{}rrrr@{}}
\toprule
$n$ & numerical & analytic & difference \\
\midrule
0 & $0.0694606784$ & $0.0694606781$ & $3\times10^{-10}$ \\
1 & $0.2008820373$ & $0.2008820344$ & $3\times10^{-9}$ \\
2 & $0.3223034063$ & $0.3223033906$ & $2\times10^{-8}$ \\
3 & $0.4337248047$ & $0.4337247468$ & $6\times10^{-8}$ \\
\bottomrule
\end{tabular}
\end{center}
The small difference is physical: the analytic formula describes the Morse
potential on the whole line, while this calculation has a wall at $x=0$.

\subsection*{Avoided crossing of two diabatic curves}

A potential defined as an eigenvalue of a $2\times2$ matrix has a closed form,
\begin{equation}
  V_{\pm}(x) = \frac{E_{1}+E_{2}}{2}
  \pm \sqrt{\left(\frac{E_{1}-E_{2}}{2}\right)^{2} + W^{2}},
\end{equation}
so it fits in a single expression. For a Morse curve $E_1$ crossing a
Coulombic one $E_2 = 1/x + \Delta$ with coupling $W = \alpha x/(x^{2}+1)$, at
$\Delta=0.5$ and $\alpha=0.3$, the lower branch is
\begin{lstlisting}[style=yaml]
physics: {mass: 100.0, hbar: 1.0}
bspline: {n_nodes: 301, order: 12, domain: [0.0, 100.0]}
potential:
  - "{'domain': '(0, 100]', 'function': '0.5*((1-exp(-(x-1)))^2 + 1/x + 0.5) - sqrt(0.25*((1-exp(-(x-1)))^2 - 1/x - 0.5)^2 + (0.3*x/(x^2+1))^2)'}"
\end{lstlisting}
whose lowest states come out at $0.0543$, $0.1852$, $0.3059$, $0.4193$,
$0.5109$, $0.5116$. The near-degenerate pair at the top is a clustering in the
density of states --- the signature of quasi-bound states behind the barrier.

To isolate such a state, multiply the tabulated wavefunction by a smooth step
$f(x)$ centred on the barrier, project the result onto the spectrum with
\texttt{coefficients()}, normalize, and propagate. Its survival probability
then shows the decay.

\subsection*{Hydrogen}
\begin{lstlisting}[style=yaml]
physics: {mass: 1.0, hbar: 1.0}
bspline: {n_nodes: 201, order: 12, domain: [0.0, 160.0]}
potential:
  # l = 0:
  - "{'domain': '(0, 160]', 'function': '-1/x'}"
  # l = 1 instead:
  # - "{'domain': '(0, 160]', 'function': '-1/x + 1/x^2'}"
\end{lstlisting}
Equation~\eqref{eq:tise} is already the reduced radial equation for
$u_{n\ell}(r)$, so the $\psi$ in the output \emph{is} $u_{n\ell}$ --- use it as
it stands, without dividing by $r$.

At $\ell=1$ the solver imposes $u(0)=0$ and the centrifugal term $1/r^{2}$
supplies the additional suppression at the origin variationally: with $R=100$
the states come out at $-0.125$, $-0.0555\overline{5}$ and $-0.03125$, with
differences of $4\times10^{-15}$, $1\times10^{-16}$ and $3\times10^{-13}$ from
$-1/2n^{2}$ for $n=2,3,4$.

% =====================================================================
\section{Potentials with steps, kinks and singularities}
\label{sec:structure}

Write the potential in whatever pieces the physics calls for; the basis adapts
at each junction on its own.

The adaptation is the one you would choose by hand. A B-spline basis of order
$k$ with a knot of multiplicity $m$ is $C^{k-1-m}$ at that knot. An ordinary
grid node has $m=1$ and so gives $C^{k-2}$, smoother than $\psi$ actually is
where $V$ jumps. The solver classifies every junction between pieces and adds
degenerate knots to match the smoothness $\psi$ really has:

\begin{center}
\begin{tabular}{@{}llccl@{}}
\toprule
Junction & What jumps & $\psi$ regularity & multiplicity & extra knots \\
\midrule
Step  & $V$ itself  & $C^{1}$ ($\psi''$ jumps)   & $k-2$ & $k-3$ \\
Kink  & $V'$ only   & $C^{2}$ ($\psi'''$ jumps)  & $k-3$ & $k-4$ \\
Singular & $V \to \infty$ & --- & $k-1$ & $k-2$; the one B-spline nonzero at the point is removed \\
\bottomrule
\end{tabular}
\end{center}

The junction is also spliced into the grid as a node. Quadrature then runs
interval by interval between consecutive nodes, 64-point Gauss--Legendre on
each, so every quadrature interval lies entirely on one side of the
discontinuity and sees a smooth integrand.

Against the exact transcendental solution for a square well, the four bound
states come out to $8\times10^{-14}$ with the step at $x=10$, which falls on
the uniform grid, and to $2\times10^{-14}$ with the step at $x=10.5$, which does
not. Put discontinuities wherever the physics puts them.

\subsection{Kinetic energy at a kink}

The kinetic matrix element is
\begin{equation}
  T_{ij} = \frac{\hbar^{2}}{2m}\int B_i'(x)\,B_j'(x)\,dx,
  \label{eq:weak}
\end{equation}
the weak (variational) form, with the integration by parts already carried out
and the boundary term vanishing under the Dirichlet conditions.

This is what lets a kink be handled exactly. Where $\psi$ has a kink, $\psi''$
contains a $\delta$-function; in~\eqref{eq:weak} only first derivatives appear,
and those are well defined for a $C^{0}$ function. The variational structure
does the rest.

% =====================================================================
\section{\texorpdfstring{$\delta$}{delta}-function potentials}
\label{sec:delta}

Declare Dirac deltas in their own section:

\begin{lstlisting}[style=yaml]
potential:                                     # the smooth part, as always
  - "{'domain': '[-20, 20]', 'function': '0'}"

potential_deltas:                              # adds strength * delta(x - x0)
  - {x: 0.0, strength: -1.0}
\end{lstlisting}

Each entry adds $\texttt{strength}\cdot\delta(x - x_0)$ to the potential, with
$x_0$ the entry's \texttt{x} field and \texttt{strength} negative for an
attractive well. Place $x_0$ strictly inside the domain.
\file{examples/delta-well.yaml} is ready to run.

\paragraph{What the solver does with them.} Two things, both automatic.

First, it sets the knot multiplicity. A $\delta$ at $x_0$ leaves $\psi$
continuous and makes $\psi'$ jump, so the basis is made exactly $C^{0}$ there:
multiplicity $k-1$, that is $k-2$ degenerate knots beyond the ordinary node,
spliced in if $x_0$ was not already a grid point. That multiplicity is what
makes the result exact:

\begin{center}
\begin{tabular}{@{}rrcrr@{}}
\toprule
extra knots & multiplicity & continuity & $E_0$ & difference \\
\midrule
0           & 1 & $C^{6}$ & $-0.42090181$       & $8\times10^{-2}$ \\
2           & 3 & $C^{4}$ & $-0.45229411$       & $5\times10^{-2}$ \\
4           & 5 & $C^{2}$ & $-0.47977270$       & $2\times10^{-2}$ \\
$6 = k{-}2$ & 7 & $C^{0}$ & $-0.50000000000002$ & $2\times10^{-14}$ \\
\bottomrule
\end{tabular}
\end{center}

\noindent($k=8$, $\lambda=1$; the last row is what the solver uses.)

Second, it adds the matrix element, which for a $\delta$ is a point evaluation
rather than a quadrature:
\begin{equation}
  \langle B_i |\, \lambda\,\delta(x-x_0) \,| B_j \rangle
  = \lambda\,B_i(x_0)\,B_j(x_0).
\end{equation}
Nothing here evaluates $\psi''$, where the $\delta$ induced by the kink would
live: the kinetic term is the weak form~\eqref{eq:weak}, so only first
derivatives appear.

\paragraph{Accuracy.} The single attractive well $V = -\lambda\delta(x)$ has
$E_0 = -\lambda^{2}/2$ and $\psi_0 = \sqrt{\lambda}\,e^{-\lambda|x|}$. On
$[-20,20]$ with \texttt{n\_nodes}${}=81$, \texttt{order}${}=8$:

\begin{center}
\begin{tabular}{@{}rrrr@{}}
\toprule
$\lambda$ & $E_0$ computed & exact & difference \\
\midrule
$0.5$ & $-0.124999998969$ & $-0.125$ & $1\times10^{-9}$ \\
$1.0$ & $-0.500000000000$ & $-0.5$   & $2\times10^{-14}$ \\
$2.0$ & $-1.999999999992$ & $-2.0$   & $8\times10^{-12}$ \\
\bottomrule
\end{tabular}
\end{center}

and $\psi(x)/\psi(0)$ follows $e^{-|x|}$ to eight digits. A weakly bound
$\delta$ spreads further --- at $\lambda = 0.5$,
$e^{-0.5\times 20} = 4.5\times10^{-5}$ is still appreciable at the wall --- so
widen the domain as $\lambda$ decreases.

\paragraph{Several at once.} Deltas accumulate, so the double well
$V = -\lambda[\delta(x+a) + \delta(x-a)]$ is two entries:
\begin{lstlisting}[style=yaml]
potential_deltas:
  - {x: -1.0, strength: -1.0}
  - {x:  1.0, strength: -1.0}
\end{lstlisting}
producing the even/odd doublet with $\kappa = \lambda(1 \pm e^{-2\kappa a})$ and
the even state lower. Deltas also combine with any smooth \texttt{potential} ---
a $\delta$ in a harmonic well, on a step, or on a Coulomb tail all work.

\paragraph{Quadrature on a cusped state.} A $\delta$ makes $\psi$ cusped, and
the trapezoid rule converges as $O(h^{2})$ across a cusp. On a 1001-point
output grid, $\int|\psi|^{2}dx$ evaluated with \texttt{np.trapezoid} reads
$1.00053$; the eigenvector itself is exactly normalized, and the difference is
the output-grid quadrature. It scales cleanly ---
$5.3\times10^{-4}$, $1.3\times10^{-4}$, $3.3\times10^{-5}$, $8.3\times10^{-6}$
as \texttt{n\_pts\_eigenstate} doubles from 1001 to 8001 --- so raise
\texttt{n\_pts\_eigenstate} to the accuracy you need.

% =====================================================================
\section{Continuum states and phase shifts}
\label{sec:continuum}

Turn them on with:
\begin{lstlisting}[style=yaml]
tise:
  continuum:
    enabled:     true
    E_threshold: 0.0     # bottom of the continuum, i.e. V at infinity
    E_max:       0.5
    n_energies:  50
    n_pts:       500
\end{lstlisting}

The solver builds scattering states and extracts $\delta(E)$ by matching to the
asymptotic form at $x_{\max}$. Choose $x_{\max}$ in a region where $V$ has
flattened out, so the matching sees the asymptotic behaviour.

For a potential with a genuine $1/x$ Coulomb tail, declare the domain unbounded
--- \texttt{'(0, inf)'} rather than \texttt{'(0, 100]'} --- so the tail is
recognized, and set \texttt{tise.continuum.l} to the angular momentum of the
$\ell(\ell+1)/2x^{2}$ term you wrote into \texttt{function}. The two must agree:
a centrifugal term decays faster than $1/x$, so the tail fit takes $\ell$ from
this field rather than reading it off the potential.

Validated on a square well of depth~1 and width~10 with a wall at the origin,
against $\tan(ka+\delta) = (k/K)\tan(Ka)$: at \texttt{n\_nodes}${}=201$ the
largest difference over 50 energies in $(0, 0.5]$ is $1.8\times10^{-12}$.

\texttt{n\_energies} also sets how many \file{continuum\_NNN.png} figures the
analysis stage writes, so 50 energies give 50 plots.

\paragraph{Energies near a box state.} In a finite box the continuum is
represented by a dense set of discrete states, and a requested energy landing on
one of them gives a state dominated by that box resonance. The solver reports
this:
\begin{lstlisting}[style=sh]
continuum energy grid point E=0.04 is within 0.00083 of confined eigenvalue
E_12=0.0408; this energy's continuum state is likely a finite-box
discretization artifact, not a physical feature
\end{lstlisting}
Shift the energy grid, or enlarge the box so the box states lie closer
together.

% =====================================================================
\section{Checking a result}

\begin{enumerate}[leftmargin=1.5em]
\item Read \file{warnings.json} and act on what it reports.
\item Double \texttt{n\_nodes}, re-run, and confirm the numbers hold.
\item Check that the highest state you use fits inside the box, from its
      classical turning point.
\item Select the bound states from \file{eigenvalues.dat} by $E < V(\infty)$.
\item For time evolution, confirm $\sum_n |c_n|^{2} = 1$ and that
      $\int|\psi(x,t)|^{2}dx$ holds at 1 across the whole time range.
\item For continuum work, converge on $\delta(E)$.
\item Compare against whatever closed form or limiting case the problem has:
      harmonic levels, the Morse formula, hydrogenic $-1/2n^{2}$, a tunneling
      splitting, a square-well phase shift, $-\lambda^{2}/2$ for a $\delta$
      well.
\end{enumerate}

\vfill
\noindent\rule{\textwidth}{0.4pt}\\
\small
Further reading in the repository:
\file{docs/guides/input-file-authoring-guide.md} for the complete input-file
reference;
\file{TISE/README.md} for the build system and the numerical methods;
\file{examples/confining-well.yaml} and \file{examples/delta-well.yaml} for
annotated starting points.

\end{document}
